\documentclass{article}
\usepackage[UTF8]{ctex}
\usepackage{amsmath,amssymb,amsfonts}
\usepackage{enumitem}
\usepackage{geometry}
\geometry{a4paper, margin=1in}
% === 错题本排版优化（不参与题目标识） ===
\setlength{\parindent}{0pt}
\setlength{\parskip}{0.35em plus 0.10em minus 0.08em}
\setlength{\emergencystretch}{3em}
\setlist{itemsep=0.25em, topsep=0.40em, parsep=0pt}
\allowdisplaybreaks[2]
% === 错题本排版优化结束 ===
\begin{document}
\section*{北京科技大学2024年数学分析真题}
\section*{11．北京科技大学2024年数学分析试卷}

\subsection{题目1}
1 ．求极限 \(\lim _{x \rightarrow 0}\left(\frac{1}{x^{2}}-\frac{1}{x \tan x}\right)\) ．

\subsection{题目2}
2．求极限 \(\lim _{x \rightarrow 0} \frac{1+\frac{1}{2} x^{2}-\sqrt{1+x^{2}}}{\left(\cos x-e^{x^{2}}\right) \sin x^{2}}\) ．

\subsection{题目3}
3．设 \(g(0)=0, g^{\prime}(0)=1\) ，分析 \(f(x)=\left\{\begin{array}{l}g(x) \sin \frac{1}{x}, \quad x>0 ; \\ g(x) \cos x, \quad x \leq 0 .\end{array}\right.\) 在 \(x=0\) 处的连续性和可导性．

\subsection{题目4}
4．求 \(I=\int x e^{a x} \cos b x \mathrm{~d} x\) 及 \(J=\int x e^{a x} \sin b x \mathrm{~d} x\) ．

二．（10 分）设正数列 \(\left\{a_{n}\right\}\) 满足 \(a_{n}=\frac{a_{n+1}^{2}}{n}+a_{n+1}, n=1,2, \cdots\) ．计算极限 \(\lim _{n \rightarrow \infty} a_{n} \ln n\) ．

三．（10 分）设 \(f(x)\) 在 \([a, b]\) 上有二阶导数，且 \(f^{(k)}(a)=f^{(k)}(b)=0(k=0,1)\) ，证明：存在 \(\xi \in(a, b)\) ，使得

\[
f^{(2)}(\xi)=f(\xi)
\]

四．（10 分）利用致密性定理证明闭区间上的连续函数必然是有界的．

五．（15 分）按照 \(p\) 的范围来说明级数 \(\sum_{n=1}^{\infty}\left(\frac{1}{n^{p}}-\ln \left(1+\frac{1}{n^{p}}\right)\right)(p>0)\) 的收敛性。

六．（15 分）求 \(x=r \cos \varphi, y=r \sin \varphi, z=r \cot \alpha\) 在点 \(M_{0}\left(r_{0}, \varphi_{0}\right)\) 处的切平面方程（ \(\alpha\) 为常数）．

七．（10 分）讨论级数 \(\sum_{n=0}^{\infty} \int_{0}^{1}(-1)^{n}(1-x) x^{n} \mathrm{~d} x\) 的收敛性并计算其和．

八．（15 分）设函数 \(f(x)\) 连续，\(\sum\) 是球面 \(x^{2}+y^{2}+z^{2}=1\) ，且 \(a, b, c\) 是常数，证明：

\[
\iint_{\sum} f(a x+b y+c z) \mathrm{d} S=2 \pi \int_{-1}^{1} f\left(\sqrt{a^{2}+b^{2}+c^{2}} u\right) \mathrm{d} u
\]

九．（15 分）设 \(a>0, b>0, m \neq 0\) 为常数，计算积分

\[
I(a, b)=\int_{0}^{+\infty} \frac{e^{-a x}-e^{-b x}}{x} \cos (m x) \mathrm{d} x
\]

十．（10 分）设 \(P(x, y), Q(x, y)\) 在 \(\mathbb{R}^{2}\) 上有一阶连续偏导数，若对以任意点 \(\left(x_{0}, y_{0}\right) \in \mathbb{R}^{2}\) 为中心，以任意\\
\(r>0\) 为半径的上半圆周 \(L_{r}: y-y_{0}=\sqrt{r^{2}-\left(x-x_{0}\right)^{2}}\) ，均有

\[
I(r)=\int_{L_{r}} P(x, y) \mathrm{d} x+Q(x, y) \mathrm{d} y=0
\]

证明：\(P(x, y) \equiv 0, \frac{\partial Q(x, y)}{\partial x} \equiv 0\) ．
\end{document}
